% Generated draft: requires mathematical and source review.
\documentclass[11pt]{article}
\usepackage[margin=25mm]{geometry}
\usepackage{fontspec}
\setmainfont{Latin Modern Roman}
\usepackage{amsmath,amssymb,mathtools,unicode-math}
\setmathfont{Latin Modern Math}
\usepackage{enumitem,needspace,xurl,hyperref}
\setlength{\parindent}{0pt}
\setlength{\parskip}{5pt}
\setlength{\emergencystretch}{3em}
\newcommand{\N}{\mathbb{N}}
\newcommand{\Z}{\mathbb{Z}}
\newcommand{\Q}{\mathbb{Q}}
\newcommand{\R}{\mathbb{R}}
\newcommand{\C}{\mathbb{C}}
\newcommand{\F}{\mathbb{F}}
\newcommand{\A}{\mathbb{A}}
\newcommand{\PP}{\mathbb{P}}
\newcommand{\E}{\mathbb{E}}
\newcommand{\eps}{\varepsilon}
\newcommand{\dd}{\,\mathrm{d}}
\begin{document}
\section*{UnsolvedMath ID: 20002911}
\textbf{Source problem code:} catalogue entry
\subsection{Definitions and mathematical statement}
The source asks for a representation-theoretic explanation of charge and a relation between filtration-adapted parametrizations and Lusztig--Kashiwara canonical bases. It does not specify a representation, weight, filtration, or meaning of ``non-miraculous.'' The aid uses the type-$A$ Brylinski--Kostant (BK) filtration and distinguishes equality of graded dimensions from a canonical filtered-basis identification.

Let $V_\lambda$ be a polynomial irreducible $\mathrm{GL}_n(\C)$ representation, $\mu$ a dominant weight of the same degree, and $e$ a chosen principal raising nilpotent. On $V=V_\lambda[\mu]$ define
\[
 F_i=V\cap\ker(e^{i+1}),\qquad F_{-1}=0.
\]
Let $E$ have formal basis indexed by semistandard tableaux $T$ of shape $\lambda$ and content $\mu$, and let $C_i$ be the span of basis elements with $\operatorname{charge}(T)\le i$. With the appropriate dominant-weight and Kostka--Foulkes conventions, the aid invokes
\[
 \sum_i\dim(F_i/F_{i-1})t^i
 =K_{\lambda\mu}(t)
 =\sum_Tt^{\operatorname{charge}(T)}.
\]
Thus the two flags have the same dimension vector, but this does not say that the canonical basis itself is adapted to $F$.

The reported exact reduction considers
\[
 \mathcal T(C,F)=\{g:E\xrightarrow{\sim}V:g(C_i)=F_i\text{ for every }i\}.
\]
It is nonempty. Let $P_C$ and $P_F$ be the parabolic stabilizers of the respective flags, omitting repeated steps. Left composition by $P_F$ and right composition by $P_C$ make it a transporter bitorsor, each action free and transitive. Existence of many such isomorphisms is automatic from the dimension equality; the missing geometric/canonical datum is a distinguished point or otherwise justified choice, not another equality of Hilbert series.

For $V_\lambda=\operatorname{Sym}^d(\C^n)$ and dominant $\mu=(\mu_1,\ldots,\mu_n)$, $\sum_j\mu_j=d$, the weight space is one-dimensional. In polynomial coordinates take $e=\sum_{j=1}^{n-1}x_j\partial_{x_{j+1}}$. The aid reports that its unique global-basis line first appears in BK degree
\[
 n(\mu)=\sum_j(j-1)\mu_j,
\]
matching the charge of the unique one-row tableau of that content. The global-basis vector is proportional to the weight monomial; its normalization is not fixed by this filtration statement. In this case the transporter ambiguity reduces to a nonzero scalar.

The aid also attributes to Patimo (2025) an affine-Grassmannian geometric/crystal explanation of type-$A$ charge not based on cyclage. It says lifting this to MV cycles and distinguished filtered bases remains open in the literature checked. Full hypotheses and definitions are absent, so neither a canonical comparison in general nor extension to other types is asserted.
\subsection{Short English statement}
Charge and the BK filtration have matching graded dimensions in the supplied type-$A$ setting. The aid identifies the extra choice needed to compare actual filtered bases and checks the one-row case exactly. Its cited geometric explanation does not yet supply the requested general canonical-basis lift.
\subsection{Sources}
\begin{itemize}
\item Local UnsolvedMath registry, permanent ID 20002911, source code catalogue entry; created and updated 14 August 2026. Canonical \texttt{aim-special-functions-notes.json}, zero-based \texttt{notes[15]}, tag \texttt{problem}, item 16.
\item AIM ``Generalized Kostka polynomials'': \url{https://aimath.org/WWN/kostka/prob2.pdf}. Wording is marked exact; original remarks and literature field are empty. Lusztig--Kashiwara bases are named without a complete bibliographic reference.
\item Research attempt 1: \texttt{partial\_result}, problem status at run \texttt{partially\_solved}, rigor label \texttt{proved}, novelty confidence low. It invokes BK/Kostka results and Patimo (2025), but supplies no full titles, theorem statements, or authoritative bibliographic links.
\item Ulam AI / UnsolvedMath Contributors, CC BY 4.0. Machine-generated research aid, not peer reviewed; independently verify before citation.
\end{itemize}
\end{document}
